\documentclass[11pt]{article}
\usepackage[T1]{fontenc}
\usepackage{lmodern}
\usepackage{amsmath,amssymb,amsthm,mathtools}
\usepackage[margin=1.05in]{geometry}
\usepackage{microtype}
\usepackage{enumitem}
\PassOptionsToPackage{hyphens}{url}
\usepackage[hidelinks]{hyperref}
\hypersetup{pdftitle={A fixed smooth metric violating Yau's nodal upper bound},
 pdfauthor={Kaiwen Jin},
 pdfsubject={An attributed exposition of OpenAI's construction on S4 times S1}}
\newcommand{\R}{\mathbb R}
\newcommand{\C}{\mathbb C}
\newcommand{\E}{\mathbb E}
\newcommand{\Prob}{\mathbb P}
\newcommand{\dd}{\mathrm d}
\newcommand{\HH}{\mathcal H}
\newcommand{\supp}{\operatorname{supp}}
\newcommand{\Div}{\operatorname{div}}
\newcommand{\tr}{\operatorname{tr}}
\newcommand{\Z}{\mathcal Z}
\newtheorem{theorem}{Theorem}[section]
\newtheorem{proposition}[theorem]{Proposition}
\newtheorem{lemma}[theorem]{Lemma}
\theoremstyle{definition}
\newtheorem{definition}[theorem]{Definition}
\theoremstyle{remark}
\newtheorem{remark}[theorem]{Remark}
\numberwithin{equation}{section}
\setlist[enumerate]{itemsep=3pt,topsep=5pt}
\emergencystretch=1em
\allowdisplaybreaks[1]
\title{A fixed smooth metric violating Yau's nodal upper bound\\
 \large An exposition of OpenAI's construction}
\author{Kaiwen Jin\\
 \small\href{mailto:yui.ui.siki@gmail.com}{yui.ui.siki@gmail.com}}
\date{October 12, 2026}
\begin{document}
\maketitle
\begin{abstract}
We present the construction of OpenAI's September 2026 preprint
\emph{Power-law violations of Yau's nodal upper bound}.
It produces one smooth Riemannian metric on \(S^4\times S^1\) and a
sequence of real exact eigenfunctions whose intrinsic nodal four-volume
exceeds the square root of the eigenvalue by a fixed positive power.
The account follows the local envelope and packet constructions through
exact coefficient correction, spectral persistence, and the smooth
limit. Particular attention is given to the common exponent and to the
order in which frequencies and perturbation tolerances are chosen.
The counterexample and its mathematical mechanism are due to OpenAI;
this note is an attributed exposition and literature update.
\end{abstract}

\section{The question and the result}\label{sec:intro}

The smooth upper-bound question asks whether, on every closed connected
smooth Riemannian manifold \((M,g)\) of dimension \(d\ge2\), there is a
finite constant \(C(M,g)\) such that
\begin{equation}\label{eq:question}
 \HH_g^{d-1}(\Z(u))\le C(M,g)\sqrt{\lambda},
 \qquad -\Delta_g u=\lambda u,\quad \lambda>0,
\end{equation}
for every nonzero real eigenfunction. Here
\(\Z(u)=\{x:u(x)=0\}\), and Hausdorff measure is taken for the distance
of the same metric \(g\). The constant must be uniform over eigenvalues
and eigenfunctions. This is the formulation of AIM Problem~011 at the
revision recorded in \cite{AIM011}.

OpenAI's preprint \cite[Theorem~1.1]{OpenAI2026} gives the following
counterexample. All substantive construction below follows that source;
the exposition does not claim a new counterexample.

\begin{theorem}[OpenAI]\label{thm:main}
There are a smooth positive Riemannian metric \(g\) on \(S^4\times S^1\),
a number \(\epsilon_0>0\), and nonzero real smooth functions \(u_j\)
satisfying
\[
 -\Delta_g u_j=\lambda_j u_j,\qquad
 \lambda_j>0,\qquad \lambda_j\longrightarrow\infty,
\]
such that
\begin{equation}\label{eq:main}
 \frac{\HH_g^4(\Z(u_j))}
      {\lambda_j^{1/2+\epsilon_0}}\longrightarrow\infty.
\end{equation}
\end{theorem}

One admissible example in dimension five negates the universal assertion
\eqref{eq:question}. Indeed, \eqref{eq:main} makes the ratios
\(\HH_g^4(\Z(u_j))/\sqrt{\lambda_j}\) unbounded for this fixed metric.
This deduction uses the quantifier over manifolds in the question; it
does not assert a separate result in every dimension.

The construction first works with a weighted elliptic operator on
\(S^4\), whose tensor and density can be perturbed independently. Local
complex waves create many pairs of disks on which a real function has
opposite signs. An explicit correction turns this function into an exact
eigenfunction with arbitrarily small changes of the coefficients. A
protected round region then permits a further perturbation that makes
its eigenvalue simple while retaining the function. Simplicity and
strict signs allow every earlier construction to survive later changes.
Summable perturbations in every finite derivative norm produce one
smooth limiting operator. A warped circle realizes that operator as the
ordinary Laplacian of a fixed five-dimensional metric.

The quantitative point is that the local gain has one exponent
\(\gamma>0\), independent of the finite accuracy subsequently requested.
At stage \(j\), the frequency can therefore absorb its stage-dependent
certificate constant. Sections~\ref{sec:envelope}--\ref{sec:exact}
prove the local construction. Section~\ref{sec:limit} supplies the
single metric and the intrinsic measure estimate.

\section{The weighted operator and a quantitative step}\label{sec:base}

Let \(B=S^4\subset\R^5\), with round metric \(g_0\) and round density
\(\mu\). Write points as \((x_1,x_2,y)\), \(y\in\R^3\), and put
\(t=|y|^2\). A positive pair \(\mathcal P=(A,\rho)\) consists of a
smooth positive definite symmetric contravariant tensor \(A\) and a
smooth positive function \(\rho\). Its principal metric and weighted
operator are
\begin{equation}\label{eq:operator}
 G^{-1}=A/\rho,\qquad
 Lf=\rho^{-1}\Div_\mu(A\,\dd f).
\end{equation}
Thus \(-L\) is nonnegative and symmetric in \(L^2(\rho\,\dd\mu)\),
with energy \(\int_B A(\dd f,\dd f)\,\dd\mu\) and form domain \(H^1(B)\).
The shifted energy obtained by adding
\(\int_B\rho f^2\,\dd\mu\) is coercive.

Choose a finite smooth atlas with compact subpatches covering \(B\).
The \(C^r\) norms include all coordinate derivatives through order \(r\),
and the pair norm is the sum of its component norms. We use monotone
norms in \(r\). All pairs stay in one small \(C^1\) neighborhood of
\((g_0^{-1},1)\); hence their metrics, inverse metrics and connections
have common bounds and common ellipticity constants. Higher derivative
bounds need only be finite for a fixed input pair.

Envelope gradients and Hessians refer to \(G\). For a real function \(h\)
define
\begin{equation}\label{eq:test}
 p=\nabla_G h,\quad Q=(1+|p|_G^2)^{1/2},\quad
 J_h(e)=\nabla_G^2h(p/Q,p/Q)+\nabla_G^2h(e,e),
 \quad |e|_G=1,\quad \dd h(e)=0.
\end{equation}
At \(p=0\) every unit vector is admissible. Complex tensor evaluations
are bilinear; norms of complex vectors and matrices are Hermitian.
Fix \(\delta=10^{-6}\). Constants \(c,C,C_r\) can change between
occurrences; their important dependencies are specified locally.

For an integer \(N>1\), set
\begin{equation}\label{eq:beam}
 k=\sqrt{N(N+3)},\qquad u_*=\operatorname{Re}(x_1+i x_2)^N.
\end{equation}
The polynomial is homogeneous and harmonic in \(\R^5\). Substitution
into
\(\Delta_{\R^5}=\partial_r^2+4r^{-1}\partial_r+r^{-2}\Delta_{g_0}\)
gives \(-\Delta_{g_0}u_*=k^2u_*\). Where \(t<1\),
\[
 u_*=(1-t)^{N/2}\cos(N\vartheta),
\]
with \(\vartheta\) the angle in the \((x_1,x_2)\)-plane.
Part of a round region will retain this exact solution.

\begin{definition}\label{def:sign}
In a fixed coordinate chart \(\Omega\subset\{t>1/2\}\), a sign
certificate for a continuous real function is a finite collection of
pairs of congruent closed parallel three-disks \(D_a^-,D_a^+\).
Their joining closed cylinders lie compactly inside pairwise disjoint
open coordinate balls. The function has strictly opposite signs on the
two disks of each pair. The certificate size is
\[
 |\mathcal C|=\sum_a\HH_{\mathrm{Euc}}^3(D_a^+).
\]
\end{definition}

Every joining segment contains a zero. Strict signs on finitely many
compact disks persist under a sufficiently small \(C^0\) change, and
under multiplication by any nonzero real scalar.

\begin{proposition}[OpenAI's quantitative step]\label{prop:step}
There are a \(C^1\) neighborhood \(\mathcal U\) of the round pair,
a chart \(\Omega\) as above, and a number \(\gamma>0\), with this
property. Fix a smooth \((A,\rho)\in\mathcal U\) that is round on
\(\{t<l_+\}\), where
\[
 1/4\le l_-<l_+\le1/2.
\]
For every integer \(r\ge0\) and every \(\eta>0\) there are \(c>0\)
and \(N_0\) such that, for every integer \(N\ge N_0\), there are a
smooth positive pair \((\widehat A,\widehat\rho)\in\mathcal U\) and
a nonzero real smooth function \(u\) for which
\begin{align}
 \|(\widehat A-A,\widehat\rho-\rho)\|_{C^r}&<\eta,
 &(\widehat A,\widehat\rho)&=(g_0^{-1},1)\quad(t<l_-),\label{eq:stepmove}\\
 -\widehat L u&=k^2u,
 &k^2&\text{ is simple on }B,\label{eq:stepexact}\\
 |\mathcal C|&\ge c k^{1+\gamma}
 &&\text{for a sign certificate in }\Omega.\label{eq:stepsigns}
\end{align}
Here \(k\) is given by \eqref{eq:beam}. The neighborhood, chart and
exponent are common to all steps. The constant \(c\) and threshold
\(N_0\) may depend on the fixed pair, collar, norm order and tolerance.
\end{proposition}

The next three sections give the proof, following
\cite[Proposition~2.2 and Sections~3--5]{OpenAI2026}.
The location of \(c\) before the universal frequency quantifier in this
statement will be essential in Section~\ref{sec:limit}.

\section{A common power gain in the envelope}\label{sec:envelope}

We construct \(h=h_k\) with a positive test \eqref{eq:test} wherever
packets will be centered, while making the integral of its local
oscillation scale grow by a power of \(k\).

\begin{proposition}\label{prop:envelope}
There are a small \(C^1\) neighborhood \(\mathcal G\) of \(g_0\),
fixed regions \(V\Subset\Omega\subset\{t>1/2\}\), and common
\(a_*,c,\gamma>0\), such that every fixed smooth \(G\in\mathcal G\)
admits envelopes \(h_k\) satisfying
\begin{equation}\label{eq:envelope}
\begin{gathered}
 h_k=-a_*t+c_k\quad(t\le1/2),\qquad
 \|h_k\|_{C^j}\le C_j k^{\delta(j+1)}\quad(j\ge0),\\
 \max_e J_{h_k}(e)\ge c k^{-\delta}\quad(t\ge1/4),\qquad
 \int_V(1+|\nabla_Gh_k|_G^2)^{1/2}\,\dd x\ge c k^\gamma .
\end{gathered}
\end{equation}
The additive constants \(c_k\) may be any prescribed bounded sequence.
For each fixed derivative order the constant and frequency threshold
can depend on that order and on the fixed smooth metric.
The exponent \(\gamma\) is independent of these orders, all later
residual accuracies, and all round gluing collars.
\end{proposition}

\subsection{A seed and a repeated local modification}

The round seed \(\widetilde h=-a_*t\) obeys
\[
 |\nabla_{g_0}\widetilde h|^2=4a_*^2t(1-t),\qquad
 \nabla_{g_0}^2\widetilde h=2a_*(tg_0-\dd y\cdot\dd y).
\]
Since \(\dd y:T_xB\to\R^3\) has a four-dimensional domain, there
is a unit vector in its kernel. It is an admissible test, with Hessian
value \(2a_*t\). The Hessian norm is at most \(2a_*\), so for
\(t\ge1/4\) and \(a_*^2\le1/8\),
\[
 J_{\widetilde h}(e)\ge a_*/2-2a_*^3\ge a_*/4.
\]
At a point with \(t=1\), the Hessian has two eigenvalues \(2a_*\)
and two zero eigenvalues. Add a sufficiently small smooth quadratic
bump there, supported in \(\{t>1/2\}\), to make the Hessian positive
definite at that point. Throughout the bump support two positive
directions remain and the gradient stays sufficiently small. Their
two-plane intersects \(\ker\dd h\), leaving a positive test even at
a critical point. Away from that support the old kernel test applies.
Choose a nearby region \(U_0\) where the Hessian is positive definite
and the gradient is separated from zero. These properties persist in
a common small \(C^1\) metric neighborhood: the Hessian of a fixed
function changes by the connection difference contracted with its
differential. Fix \(U_0\Subset V\Subset\Omega\).

A cube \(P\) of side \(\nu S\), \(S\le1\), is called prepared for \(h\)
if, on its closure,
\begin{equation}\label{eq:prepared}
\begin{gathered}
 |\nabla h|>0,\qquad
 |\nabla^2h|\le C_{\rm p}|\nabla h|/S,\\
 \text{there is a two-plane }E\subset\ker\dd h
 \text{ with }J_h(e)\ge c_{\rm p}|\nabla h|/S
 \text{ for every unit }e\in E.
\end{gathered}
\end{equation}
The plane need not depend continuously on the point. The seed is
prepared at scale \(1\) on a sufficiently small cube in \(U_0\),
for fixed large \(C_{\rm p}\) and small \(c_{\rm p}>0\).

\begin{lemma}\label{lem:corrugation}
One can fix a reciprocal integer \(\nu>0\), constants
\(0<\alpha\le1\le\beta\), and \(\zeta>0\), uniformly under the
common metric bounds, so that every prepared parent \(P\) admits a
smooth modification \(v=h+w\), supported in its interior, with
\[
 \alpha|\nabla h|\le|\nabla v|\le\beta|\nabla h|,\qquad
 \max_eJ_v(e)\ge c|\nabla h|/S.
\]
Put \(\ell=\nu^2S\). Some cubes of side \(\nu\ell\) inside \(P\)
are prepared at scale \(\ell\), with the same constants, and their
union \(\mathcal R_P\) has
\begin{equation}\label{eq:gain}
 \int_{\mathcal R_P}|\nabla v|\,\dd x
 \ge(1+\zeta)\int_P|\nabla h|\,\dd x.
\end{equation}
If \(s=|\nabla h(x_P)|\) at the parent center, then
\begin{equation}\label{eq:increment}
 \|D^jw\|_\infty\le C_j s\ell^{1-j}\qquad(j\ge0).
\end{equation}
Neither the modification nor these bounds use higher derivatives
of the input \(h\).
\end{lemma}

\begin{proof}
Choose a smooth \(F:[0,1]\to\R\), positive inside, equal to its
argument near zero and flat at one. In a finite collection of
disjoint balls of radii \(R_i\) in a unit three-cube, define
\[
 \psi(r)=-\int_r^{R_i}F(q/R_i)\,\dd q,\qquad
 f=F(r/R_i),\qquad u_0=f/r.
\]
Extend by zero and periodically. At the center \(\psi\) is a constant
plus \(r^2/(2R_i)\), so it is smooth there; flatness makes the boundary
extension smooth. With radial unit vector \(\sigma\),
\begin{equation}\label{eq:radial}
 D\psi=f\sigma,\qquad
 D^2\psi=f'\sigma\otimes\sigma+u_0(I-\sigma\otimes\sigma).
\end{equation}
The inequalities \(f^2|f'|\le C u_0\) and \(f^2\le C u_0\) hold.
The first constant does not depend on a ball radius, since
\(f^2|f'|/u_0=vF(v)|F'(v)|\), \(v=r/R_i\).
Also \(u_0=0\) implies \(f=f'=\psi=0\).
At a center the radial Hessian is independent of the choice of \(\sigma\).

Fix a small amplitude \(a>0\) using that radius-independent bound.
The mean gradient factor inside any ball is
\[
 m(a)=3\int_0^1v^2\sqrt{1+a^2F(v)^2}\,\dd v>1.
\]
Let \(\eta_0=(m(a)-1)/8\). Choose the finite packing to have volume
fraction \(\varpi\) with \(\varpi m(a)>1+6\eta_0\). Such finite
packings can approach full volume: finitely many grid cubes capture
a fixed fraction of any remaining open volume; smaller inscribed
balls capture a fixed fraction of those cubes. Repeating this
finitely often leaves an arbitrarily small uncovered volume and
keeps closures disjoint. Reduce each ball to a concentric core
of relative radius \(1-\eta_{\rm c}\), so that the integral per
period over the cores is \(C_{\rm core}>1+5\eta_0\).
Choose an inner-parent volume fraction \(q_{\rm in}<1\) with
\(q_{\rm in}C_{\rm core}>1+4\eta_0\).
A smooth parent cutoff \(\chi\), equal to one on that inner cube
and flat on its zero set, is now fixed.

Take affine coordinates \(z=(z_1,z')\) that make the metric
Euclidean at \(x_P\) and \(\nabla h(x_P)=s\partial_{z_1}\).
The preparation bound and the common connection bound give
\(|D(\dd h)|\le C|\dd h|/S\). Integration along parent segments
of length \(O(\nu S)\) gives
\begin{equation}\label{eq:parent}
 \dd h/s=\dd z_1+O(\nu),\qquad G_{ij}=I+O(\nu),
 \qquad |\nabla h|/s=1+O(\nu).
\end{equation}
Rescale \(\chi\) to the parent and set
\begin{equation}\label{eq:modify}
 v=h+s\ell a\chi\psi(z'/\ell).
\end{equation}
The ratio between the pattern scale and cutoff scale is \(\nu\);
hence
\(\dd v/s=\dd z_1+a\chi f\sigma\cdot\dd z'+O(\nu)\).
This proves the gradient comparison.

We give the Hessian check also where the favorable coefficient
vanishes. Put \(D=\nabla v/\sqrt{1+|\nabla v|^2}\) and
\[
 \mathcal F=\{e:\dd v(e)=0,\ \sigma\cdot e'=0\}.
\]
This is a two-plane. For metric-unit \(e\in\mathcal F\),
\[
 |e'|^2\ge1/2,\qquad |\sigma\cdot D'|\le C_0a\chi f+C\nu.
\]
After dividing the new Hessian test by \(s/\ell\), its principal
contribution from \eqref{eq:radial} is
\[
 a\chi\{f'(\sigma\cdot D')^2+
            u_0(|\Pi_{\sigma^\perp}D'|^2+|e'|^2)\}.
\]
Choose \(a\) so that
\(2C_0^2a^2\sup_v vF(v)|F'(v)|\le1/4\).
The displayed expression is then at least
\(a\chi u_0/4-Ca\nu^2\chi|f'|\).
Mixed cutoff terms on \(e,e\) vanish. On \(D,D\) they are bounded by
\[
 C\nu f|D_{\rm sc}\chi|(\chi f+\nu).
\]
The remaining cutoff and connection errors are bounded by
\[
 C\nu^2|\psi||D_{\rm sc}^2\chi|
       +C\ell(\chi f+\nu).
\]
Here \(D_{\rm sc}\) denotes differentiation at the parent cutoff scale.

Let \(X=\chi u_0\). On the compact set of cutoff and period positions
and bounded affine frames, every continuous \(b\ge0\) that vanishes
when \(X=0\) satisfies
\[
 \sup\frac{\nu^2b}{X+\nu^2}\longrightarrow0.
\]
For \(X\ge\tau\), bound the ratio by \(\nu^2\sup b/\tau\);
for \(X<\tau\), bound it by \(\sup_{X<\tau}b\), then let
\(\tau\downarrow0\). Apply this to
\(\chi|f'|,\ f|D_{\rm sc}\chi|,\ |\psi||D_{\rm sc}^2\chi|\);
their vanishing follows from flatness and \(u_0=0\Rightarrow f=f'=\psi=0\).
The other mixed term is \(O(\nu X)\), and
\(\chi f\le C\sqrt X\), \(\ell\le\nu^2\) imply
\(\ell(\chi f+\nu)/(X+\nu^2)=O(\nu)\).
Consequently, uniformly over admissible parents,
\begin{equation}\label{eq:newhessian}
 \frac{\ell}{s}\{\nabla^2(v-h)(D,D)+\nabla^2(v-h)(e,e)\}
 \ge c_aX-\varepsilon(\nu)(X+\nu^2),\qquad
 \varepsilon(\nu)\to0.
\end{equation}

The old Hessian is \(O(\nu^2)\) in these normalized units.
If \(X/\nu^2\to\infty\), the new positive term dominates it.
If this ratio stays bounded, the old and new gradients differ by a
relative \(O(\nu)\). The map \(T(p)=p/\sqrt{1+|p|^2}\) obeys
\[
 |T(q)-T(p)|\le2\epsilon
 \quad\text{when }|q-p|\le\epsilon|p|,\quad\epsilon\le1/2.
\]
This follows by integrating its derivative along a segment whose
distance from zero is at least \(|p|/2\). A rotation \(R=I+O(\nu)\)
takes \(\ker\dd h\) to \(\ker\dd v\). Inside this three-dimensional
space, \(RE\) and \(\mathcal F\) have intersection dimension at least
\(2+2-3=1\). In that direction the old positive test changes by
\(o(\nu^2)\); the new contribution in \eqref{eq:newhessian} is
bounded below by \(-o(\nu^2)\). Thus a positive margin of size
\(c_{\rm p}\nu^2/8\) remains. A contrary sequence of configurations
has a subsequence with bounded ratio or ratio tending to infinity;
both cases exclude its failure. This proves the uniform one-direction
margin on the whole parent. The intersection calculation is the use
of base dimension four.

Where \(\chi=1\) and the period position lies in a core, \(u_0\)
has a fixed positive lower bound. Both directions in \(\mathcal F\)
then have test at least \(c's/\ell\), while the Hessian norm is at most
\((C'+O(C_{\rm p}\nu^2))s/\ell\). Fix sufficiently large
\(C_{\rm p}\) and small \(c_{\rm p}\), then sufficiently small \(\nu\);
this region is prepared at scale \(\ell\).
Retain only whole child cubes contained in it.

To compare mass, tile the affine coordinates by period cells of
side \(\ell\), including a dummy \(z_1\) period. Cells meeting the
inner-parent boundary occupy relative volume \(O(\nu)\).
Discarding child cubes meeting core boundaries loses another
\(O(\nu)\) fraction, because the finite cores have finite boundary
area in a period. On the retained region
\[
 |\nabla v|/s=\sqrt{1+a^2f^2}+O(\nu).
\]
The fixed affine Jacobian cancels against parent volume. Thus
\[
 \frac1{s|P|}\int_{\mathcal R_P}|\nabla v|\,\dd x
 \ge q_{\rm in}C_{\rm core}-C\nu,\qquad
 \frac1{s|P|}\int_P|\nabla h|\,\dd x\le1+C\nu.
\]
Taking both errors below \(\eta_0\) gives a ratio exceeding
\((1+3\eta_0)/(1+\eta_0)>1+\zeta\), with
\(\zeta=\eta_0/(1+\eta_0)>0\).

Finally, \(j\) derivatives of \eqref{eq:modify}, with \(q\) on the
cutoff, cost
\[
 s\ell(\nu S)^{-q}\ell^{-(j-q)}
   =s\ell^{1-j}\nu^q\le s\ell^{1-j}.
\]
The slope and affine frame are frozen constants in this expression.
This proves \eqref{eq:increment}. The order of choices is the
pattern and its small amplitude, the packing and cutoffs,
the preparation constants, and finally \(\nu\).
\end{proof}

\subsection{Iteration with every finite derivative order}

\begin{proof}[Proof of Proposition~\ref{prop:envelope}]
Start from one seed cube of side \(\nu\), prepared at scale \(1\),
compactly inside \(U_0\). Write \(h_0\) for the seed and \(P_0\)
for this cube. Fix \(s_{\min}>0\) with
\( |\nabla_Gh_0|_G\ge s_{\min}\) on \(P_0\), common over
\(G\in\mathcal G\); this follows from the gradient separation on
\(U_0\) and the common metric bounds. Put \(r_{\rm c}=\nu^2\).
At generation \(i\), apply Lemma~\ref{lem:corrugation} on each
selected parent at scale \(S_i=r_{\rm c}^{i-1}\).
The selected children have side \(\nu r_{\rm c}^{i}\).
Supports lie strictly inside parents, so every finite generation
is smooth. The total selected gradient mass \(I_i\) satisfies
\[
 I_i\ge(1+\zeta)I_{i-1},\qquad
 \int_V\sqrt{1+|\nabla h_i|^2}\,\dd x\ge I_i.
\]
The initial \(I_0\) has a common positive lower bound.

Gradients grow by at most \(\beta\) each generation and retain
at least a factor \(\alpha\) on the initial cube. Disjoint supports
within a generation and \eqref{eq:increment} give
\[
 \|D^j(h_i-h_{i-1})\|_\infty
 \le C_j\beta^{i-1}r_{\rm c}^{i(1-j)}.
\]
There is no factor counting cubes. The increments are added, so
their higher derivative constants are not multiplied between
generations. Summing yields one \(C_*\), independent of \(j\), with
\begin{equation}\label{eq:iteration}
 \|h_m\|_{C^j}\le C_j e^{C_*(j+1)m}.
\end{equation}
For example, take \(C_*\) larger than
\(\log\beta+|\log r_{\rm c}|+1\).
At a point last changed at generation \(i\le m\), the retained
Hessian margin is at least
\[
 c\,s_{\min}\alpha^{i-1}/r_{\rm c}^{i-1}
 \ge c\,s_{\min}\alpha^m.
\]
Later modifications vanish there, with all jets vanishing at
support boundaries. Outside the initial cube the seed is unchanged.
Enlarge \(C_*\) to obtain \(\max_eJ_{h_m}(e)\ge c e^{-C_*m}\)
on \(t\ge1/4\).

Choose once \(\theta>0\) with \(C_*\theta\le\delta\), and put
\[
 m=\lfloor\theta\log k\rfloor,\qquad
 \gamma=\theta\log(1+\zeta)>0,\qquad h_k=h_m+c_k.
\]
The derivative and Hessian bounds follow from \eqref{eq:iteration},
and the mass bound follows from
\((1+\zeta)^m\ge(1+\zeta)^{-1}k^\gamma\).
All changes lie in the seed cube above \(t=1/2\).
The constants determining \(C_*,\theta,\gamma\) use the fixed
pattern and common \(C^1\) metric bounds only. A later smooth coordinate
change contributes finite order-dependent chain-rule constants and
lower derivatives of \(h_k\), without changing these exponents.
\end{proof}

\section{Local complex waves below the envelope}\label{sec:packets}

Fix a smooth positive input pair and an envelope from
Proposition~\ref{prop:envelope}. A wave packet will have amplitude
comparable to \(e^{kh}\) at its center, decay rapidly relative to
that envelope away from the center, and have a residual as small
as any prescribed inverse power of \(k\). The finite construction
must retain the same derivative and spanning exponents when its
accuracy is increased.

\subsection{Characteristic Hessians and polynomial equations}

In \(G\)-normal coordinates at a center, let
\(p=\nabla h\), \(Q=\sqrt{1+|p|^2}\) and \(H=\nabla^2h\).
For a unit \(v\perp p\), the bilinear square of \(P=p+iQv\) is
\(-1\). Require the phase to satisfy the bilinear eikonal equation
\(G^{-1}(\dd\phi,\dd\phi)=-1\) to first order at the center.
A phase with differential \(P\) must then have Hessian \(S\)
satisfying \(SP=0\): differentiating the equation gives this condition
because the first metric derivatives vanish at the center.
The real Hessian gap \(H-\operatorname{Re}S>0\)
will keep its exponential below \(e^{kh}\).

\begin{lemma}\label{lem:hessian}
Let \(e_1,v\) be orthonormal, \(0\le b<1\), and \(H\) real
symmetric, with
\[
 |H|\le k^{4\delta},\qquad
 b^2H(e_1,e_1)+H(v,v)\ge k^{-2\delta}.
\]
For all sufficiently large \(k\), uniformly under these hypotheses,
there are real symmetric \(U,D\) with
\begin{equation}\label{eq:hessian}
 (U+iD)(be_1+iv)=0,\quad
 H-U\ge k^{-3\delta}I,\quad |U+iD|\le k^{20\delta}.
\end{equation}
\end{lemma}

\begin{proof}
For \(b>0\) the first identity is
\(Dv=bUe_1,\ De_1=-Uv/b\). The prescribed columns are compatible
with symmetry exactly when \(b^2U_{11}+U_{vv}=0\).
If \(b\ge k^{-12\delta}\), put
\[
 c_0=\frac{b^2H_{11}+H_{vv}}{1+b^2},\qquad U=H-c_0I.
\]
Relative to \(e_1,v,\{e_1,v\}^{\perp}\), a symmetric completion is
\[
 D=\begin{pmatrix}
 -u_{12}/b&bu_{11}&-u_{2E}/b\\
 bu_{11}&bu_{12}&bu_{1E}\\
 -u_{E2}/b&bu_{E1}&0
 \end{pmatrix}.
\]
Here the \(u\)'s are the corresponding blocks of \(U\).
Then \(H-U\ge\tfrac12k^{-2\delta}I\) and
\(|D|\le Ck^{16\delta}\).

For \(0\le b<k^{-12\delta}\), set \(K=k^{14\delta}\) and
\[
 U=-K(I-vv^{\mathsf t})+b^2Kvv^{\mathsf t},\qquad
 D=-bK(e_1v^{\mathsf t}+ve_1^{\mathsf t}).
\]
The identity holds also at \(b=0\). In
\(H-U-k^{-3\delta}I\), the \(v^\perp\) block is at least
\(KI/2\), while the scalar \(v\) block is at least
\[
 k^{-2\delta}-k^{-20\delta}-k^{-10\delta}-k^{-3\delta}.
\]
Its Schur complement subtracts at most \(C|H|^2/K\le Ck^{-6\delta}\).
All negative terms are smaller than \(k^{-2\delta}\) for large \(k\).
The size bounds follow from the explicit matrices.
\end{proof}

Only the Hermitian size of a complex direction is needed to solve
the subsequent polynomial equations. To see this, let \(\mathcal P_m\)
be the homogeneous complex polynomials of degree \(m\), and suppose
\(c\le|q|\le C\). Choose \(|q_i|\ge c/2\), write
\(q\cdot\partial=q_i\partial_i+D'\), and let \(I_i\) integrate
in \(X_i\) with zero constant. The operator
\(\mathcal K=D'I_i/q_i\) lowers the degree in the other variables,
so \(\mathcal K^{m+1}=0\) on \(\mathcal P_m\). Thus
\begin{equation}\label{eq:polyinverse}
 R_{q,m}=\frac{I_i}{q_i}\sum_{r=0}^m(-\mathcal K)^r,\qquad
 (q\cdot\partial)R_{q,m}=I .
\end{equation}
Its coefficient norm is bounded in terms of \(m,c,C\).
There is no division by the possibly small bilinear square \(q\cdot q\).

\subsection{Four packets with arbitrary finite accuracy}

For packets centered at \(z\), define their normalized first jets
in that center's fixed normal coordinates by
\[
 {\bf j}_z w(x)=e^{-kh(x)}
       \left(w(x),\frac{\partial_\xi w(x)}{kQ_z}\right),
 \qquad Q_z=\sqrt{1+|\nabla h(z)|^2}.
\]
The real and imaginary columns of these jets form a real matrix.
Its least spanning singular value means the square root of the
least eigenvalue of the matrix times its transpose.

\begin{proposition}\label{prop:packets}
Given an integer \(J\ge0\) and \(M>0\), finite construction orders
can be fixed so that, for every sufficiently large \(k\), every
center \(z\) with \(t(z)\ge1/4\) has four smooth complex packets.
They have support radius at most \(Ck^{-1/3-100\delta}\), and
\begin{align}
 |D_x^j w|&\le C_j k^{2j}e^{kh}\qquad(j\ge0),\label{eq:packetder}\\
 |D_x^j(L+k^2)w|&\le Ck^{-M}e^{kh}
                    \qquad(0\le j\le J).\label{eq:packetres}
\end{align}
The family spans real first jets with singular value at least
\(ck^{-3/4}\) at \(z\), and at least \(k^{-1}\) at distance
at most \(k^{-10}\) from \(z\). The support, derivative and
spanning exponents do not depend on \(J,M\). Constants and
frequency thresholds may depend on the fixed pair and chosen orders.
\end{proposition}

\begin{proof}
The envelope bounds give \(Q\le k^{3\delta}\) and
\(|H|\le k^{4\delta}\) for large \(k\).
If \(p\ne0\), take \(e_1=p/|p|\); choose a favorable test \(e_2\)
in \eqref{eq:envelope}. At \(p=0\), choose \(e_2\) first and
then any perpendicular \(e_1\). Complete an orthonormal basis and put
\[
 \eta_{\rm v}=k^{-9\delta},\quad v_2=e_2,\quad
 v_3=\frac{e_2+\eta_{\rm v}e_3}{\sqrt{1+\eta_{\rm v}^2}},\quad
 v_4=\frac{e_2+\eta_{\rm v}e_4}{\sqrt{1+\eta_{\rm v}^2}}.
\]
Use the four pairs
\[
 (p',v)=(p,v_2),(p,v_3),(p,v_4),
             (p+Qk^{-3/4}e_1,v_2).
\]
For each set \(Q'=\sqrt{1+|p'|^2}\), \(P=p'+iQ'v\),
\(b=|p'|/Q'\). Then \(P\cdot P=-1\) and
\(1\le Q'/Q\le1+k^{-3/4}\).
The change in the favorable Hessian test is at most
\(Ck^{-5\delta}+Ck^{-3/4+4\delta}\).
For the second term, differentiate \(s^2/(1+s^2)\) as
\(s=|p|\) increases by \(Qk^{-3/4}\).
These errors are smaller than the original \(ck^{-\delta}\)
margin. Lemma~\ref{lem:hessian} therefore gives \(U,D\) with
\((U+iD)P=0\) for each pair.

Each packet uses one frozen normal chart
\(\Psi_{z,E}(\xi)=\exp_z(\sum_i\xi_iE_i)\).
For a fixed smooth input metric, all fixed-order derivatives of
these charts and their inverses are uniformly bounded over centers
and unit frames on sufficiently small common balls. This follows
by solving the smooth geodesic equation on a common bounded set
and differentiating it in initial conditions, then using compactness
of the frame bundle. No spatial derivative of the frame-selection
rule is taken.

Introduce the scaled variables
\begin{equation}\label{eq:scaling}
 r_0=k^{-100\delta},\quad X=\xi/r_0,\quad
 \kappa=kr_0Q,\quad
 h_{\rm sc}(X)=\frac{h(\Psi_{z,E}(r_0X))-h(z)}{r_0Q}.
\end{equation}
For \(q\ge2\), the chain rule gives
\[
 |D_X^q h_{\rm sc}|\le C_q Q^{-1}k^{\delta(101-99q)}.
\]
These are bounded for every fixed order. If
\(L=g^{ij}(\xi)\partial_{\xi_i\xi_j}
       +\beta^i(\xi)\partial_{\xi_i}\), multiplying \(L+k^2\)
by \(r_0^2\) gives
\[
 \mathcal L+\kappa^2/Q^2,\qquad
 \mathcal L=g^{ij}(r_0X)\partial_{X_iX_j}
                   +r_0\beta^i(r_0X)\partial_{X_i}.
\]
Its fixed-order coefficient derivatives are bounded, with
\(g^{ij}(0)=\delta_{ij}\) and \(\partial g^{ij}(0)=0\).

For an integer \(T\ge3\), construct a polynomial phase with
\[
 \Phi(0)=0,\quad \partial\Phi(0)=P/Q,\quad
 \partial^2\Phi(0)=r_0(U+iD)/Q
\]
so that
\(E_\Phi=g^{ij}(r_0X)\partial_i\Phi\,\partial_j\Phi+Q^{-2}\)
vanishes to order \(T\) at zero. The constant and linear equations
are already satisfied. At homogeneous degree \(m\ge2\), the new
term \(\Phi_{m+1}\) enters as
\(2(P/Q)\cdot\partial\Phi_{m+1}\); everything else is known.
Its direction has Hermitian norm between \(1\) and \(3\), so
\eqref{eq:polyinverse} solves every equation.
For fixed \(T\) all coefficients are uniformly bounded.

Put
\[
 \mathcal T a=2g^{ij}(r_0X)(\partial_i\Phi)(\partial_j a)
                          +(\mathcal L\Phi)a.
\]
For a finite integer \(s\ge0\), solve by homogeneous degrees for
polynomials \(a_0,\ldots,a_s\), with \(a_0(0)=1\), \(a_j(0)=0\)
for \(j\ge1\), such that
\[
 \mathcal T a_0,\qquad
 \mathcal T a_j+\mathcal L a_{j-1}\quad(1\le j\le s)
\]
vanish to order \(T\).
Again the new degree-\((m+1)\) term enters via
\(2(P/Q)\cdot\partial\). Complete all of \(a_{j-1}\) before
starting \(a_j\), so the second derivatives in its right-hand side
are already determined. Every coefficient is bounded for fixed \(T,s\).

Choose a real smooth cutoff \(\omega\), equal to one on the unit
ball and supported in the radius-two ball, and define
\begin{equation}\label{eq:packet}
 a=\sum_{j=0}^s\kappa^{-j}a_j,\qquad
 w(\Psi_{z,E}(r_0X))=
 e^{kh(z)}e^{\kappa\Phi(X)}a(X)\omega(k^{1/3}X).
\end{equation}
It extends smoothly by zero and has the stated support radius.
On its support \(a=1+O(k^{-1/3})\), hence \(|a|\ge1/2\).

The real phase lies below the envelope to quadratic order:
\begin{equation}\label{eq:localization}
 \operatorname{Re}\Phi-h_{\rm sc}
 \le k^{-3/4}|X|-c k^{-106\delta}|X|^2+C|X|^3 .
\end{equation}
The linear term is zero for the first three packets. On
\(|X|\le2k^{-1/3}\), the cubic term is absorbed into half the
quadratic deficit. The remaining maximum is
\(O(k^{-3/2+106\delta})\); since
\(k^{1-100\delta}\le\kappa\le k\), the exponential in
\eqref{eq:packet} is at most \(2e^{kh}\).
On the annulus \(\tfrac14k^{-1/3}\le|X|\le2k^{-1/3}\)
it is at most
\begin{equation}\label{eq:tail}
 e^{-2k^{1/4}}e^{kh}.
\end{equation}
Indeed, the negative exponent has size at least
\(ck^{1/3-206\delta}\), and the positive linear exponent is
at most \(Ck^{-1/12}\).

Direct differentiation before the cutoff gives
\[
\begin{split}
 e^{-\kappa\Phi}(\mathcal L+\kappa^2/Q^2)(e^{\kappa\Phi}a)
 ={}&\kappa^2E_\Phi a+\kappa\mathcal T a_0\\
 &+\sum_{j=1}^s\kappa^{1-j}
       (\mathcal T a_j+\mathcal L a_{j-1})
       +\kappa^{-s}\mathcal L a_s .
\end{split}
\]
Taylor remainders and their \(q\le J\le T\) derivatives cost
\(|X|^{T-q}\). In ordinary coordinates a derivative of such a
remainder costs at most \(k^{1/3+100\delta}\);
derivatives of the exponential, amplitude and cutoff cost at
most \(Ck^2\) each. Undoing the operator scaling gives the bound
\begin{equation}\label{eq:accuracy}
 Ck^{200\delta+2J}
 \left(k^{2-T/3}+k^{1-T/3}+k^{-(1-100\delta)s}\right)e^{kh}.
\end{equation}
All cutoff commutators have polynomial factors times \eqref{eq:tail}.
Choose, before \(k\),
\[
 T>3(M+200\delta+2J+3),\qquad
 s>\frac{M+200\delta+2J+1}{1-100\delta}.
\]
Increasing the frequency threshold then proves \eqref{eq:packetres}.
The same calculation without the operator factor proves
\eqref{eq:packetder}, with an exponent independent of \(T,s\).

At the center,
\({\bf j}_z w=(1,P/Q)+O(\kappa^{-1})\).
The real leading columns contain
\((1,(|p|/Q)e_1)\) and
\((1,(|p|/Q+k^{-3/4})e_1)\); they span the value--\(e_1\)
plane with singular value at least \(ck^{-3/4}\).
The imaginary columns span \(e_1^\perp\) with singular value
at least \(ck^{-9\delta}\).
The amplitude error is \(O(k^{-1+100\delta})=o(k^{-3/4})\).
Through two derivatives \eqref{eq:packetder} also gives
\(|D_x{\bf j}_z w|\le Ck^3\). Moving by \(k^{-10}\)
changes the matrix by only \(O(k^{-7})\).
Fixed coordinate changes have bounded inverses, so the nearby
lower bound is \(k^{-1}\) for large \(k\).
Every step includes \(p=0\).
\end{proof}

\subsection{Estimates on wavelength balls}

The next estimates control all packets that meet a small ball,
including packets centered outside the region where signs will be
measured. This is needed because a value-normalized Gaussian sum
must have bounded local derivative moments.

\begin{lemma}\label{lem:scaled}
Fix the packet orders. Let \(\mathcal W\) contain all four packets
at each of finitely many centers, with
\(|\mathcal W|\le Ck^{K_1}\) for fixed \(K_1\).
Suppose every \(x\in V\) has a center within \(C_0k^{-10}\).
Define, in the fixed chart \(\Omega\),
\[
 R_x=(kQ_x)^{-1},\quad
 S_x^2=\sum_{w\in\mathcal W}|w(x)|^2,\quad
 w_x(Y)=w(x+R_xY),\quad |Y|\le2.
\]
For packets whose support meets this ball, let \(X_z\) denote
their own centered scaled coordinates. Call them non-tail when
\(|X_z(x)|<\tfrac12k^{-1/3}\), and tail otherwise.
For every fixed \(j\ge0\),
\begin{align}
 S_x&\ge ce^{kh(x)},\label{eq:variance}\\
 \sup_{|Y|\le2}|D_Y^jw_x(Y)|
 &\le C_j|w(x)| &&\text{for non-tails},\label{eq:nontail}\\
 \sup_{|Y|\le2}|D_Y^jw_x(Y)|
 &\le C_j e^{-k^{1/4}}e^{kh(x)}
                         &&\text{for tails}.\label{eq:abstail}
\end{align}
Every non-tail has nonzero center value and cutoff identically
one on the ball, and
\begin{equation}\label{eq:imaginary}
 \left|\operatorname{Im}\frac{\dd_Yw_x}{w_x}(0)\right|\ge c.
\end{equation}
Consequently,
\begin{equation}\label{eq:squares}
 S_x^{-2}\sum_w\sup_{|Y|\le2}|D_Y^jw_x|^2\le C_j,\qquad
 S_x^{-2}\sum_{\text{non-tail }w}|w(x)|^2=1-o(1).
\end{equation}
Constants and thresholds may depend on the fixed pair, orders
and \(K_1,C_0\), but not on \(k,x\).
\end{lemma}

\begin{proof}
The envelope gives \(|D Q|\le Ck^{4\delta}\).
Any intersecting packet center has
\[
 |z-x|\le Ck^{-1/3-100\delta}+C/k,\qquad
 |Q_z-Q_x|\le Ck^{-1/3-96\delta}+Ck^{-1+4\delta}=o(1).
\]
All such centers and balls lie in a common compact subchart for
large \(k\). Since \(Q\ge1\), the ratios \(Q_z/Q_x\) and the
corresponding ratios throughout the ball are bounded above and
below. From \(|\dd h|\le C Q\),
\[
 k|h(x+R_xY)-h(x)|\le C \quad(|Y|\le2).
\]
The nearest center has \(|X_z(x)|=O(k^{-10+100\delta})\).
Its center value and nearby jet estimates yield
\(|w(x)|\ge ce^{kh(x)}\) for one of its packets, proving
\eqref{eq:variance}.

With \(\kappa_z=kr_0Q_z\), fixed normal-chart bounds give
\[
 |D_Y^qX_z|\le C_qr_0^{-1}R_x^q,\qquad
 \kappa_z|D_Y^qX_z|\le
 C_q(Q_z/Q_x)(kQ_x)^{1-q}\le C_q\quad(q\ge1).
\]
For \(q=1\),
\(\kappa_z D_YX_z=(Q_z/Q_x)D_x\xi_z\) has a uniformly
bounded inverse. The displacement of \(X_z\) on the ball is
\(O(k^{-1+100\delta}/Q_x)=o(k^{-1/3})\).

For a tail, any point where the cutoff or its derivatives
contribute satisfies
\(\tfrac14k^{-1/3}\le|X_z|\le2k^{-1/3}\).
Use \eqref{eq:tail}. Each fixed positive \(Y\)-derivative of
\(\kappa_z\Phi(X_z)\) is bounded by the preceding estimates;
products with \(m\) derivatives of \(X_z\) leave
\(\kappa_z^{1-m}\le1\).
Amplitude derivatives are bounded, and
\(k^{1/3}/\kappa_z\le k^{-2/3+100\delta}\le1\)
controls cutoff derivatives. The exponential estimate therefore
gives \eqref{eq:abstail}, even when \(w(x)=0\).

For a non-tail the cutoff stays one on the ball and the amplitude
is bounded away from zero. The same derivative bounds control
the logarithmic derivative and all higher derivatives relative
to \(|w_x|\). Integrating its real logarithmic derivative along
straight segments gives \(|w_x(Y)|\le C|w(x)|\), proving
\eqref{eq:nontail}.
Moreover,
\[
 \operatorname{Im}\partial_X\Phi(X_z(x))
       =(Q'/Q_z)v+O(k^{-1/3}).
\]
Its leading norm is at least one, so the inverse bound on
\(\kappa_zD_YX_z\) gives a positive imaginary \(Y\)-gradient.
The amplitude's logarithmic derivative is only
\(O(k^{-1+100\delta}/Q_x)=o(1)\), which proves
\eqref{eq:imaginary}. Squaring the tail bounds and using
\eqref{eq:variance} bounds their total normalized contribution
by \(C_j k^{K_1}e^{-2k^{1/4}}=o(1)\).
The non-tail contributions are controlled by their center
values. The definition of \(S_x\) then proves \eqref{eq:squares}.
\end{proof}

\section{An exact simple eigenfunction after a small perturbation}
\label{sec:exact}

Fix the input pair, \(l_-,l_+\), and put
\(w_{\rm c}=l_+-l_->0\),
\(s_i=l_-+i w_{\rm c}/5\) for \(1\le i\le4\).
For the allowed \(k=\sqrt{N(N+3)}\), let
\[
 b_N(t)=\frac{N}{2k}\log(1-t).
\]
Choose the additive constant in the envelope so that
\(h(s_2)=b_N(s_2)\). It is bounded uniformly in \(N\).
For the fixed small \(a_*\), and all sufficiently large \(N\),
\[
 \partial_t(h-b_N)
       =-a_*+\frac{N}{2k(1-t)}\ge c_0>0
       \quad(0\le t\le1/2).
\]
It follows that
\begin{equation}\label{eq:gluegaps}
 h\le b_N-c w_{\rm c}\ (t\le s_1),\qquad
 h\ge b_N+c w_{\rm c}\ (s_3\le t\le s_4).
\end{equation}
Take \(0\le\beta\le1\), equal to one for \(t\le s_3\) and
zero for \(t\ge s_4\). The smooth positive weight
\begin{equation}\label{eq:weight}
 W=e^{kh}+\beta(t)e^{kb_N(t)}
\end{equation}
uses a zero extension of the second summand.
The pair is round on the support of \(\beta\).
Thus \((L+k^2)(\beta u_*)\) is only a cutoff commutator.
Its support lies in the exponential gap of \eqref{eq:gluegaps}.
For every fixed \(q\ge0\) and \(M>0\),
\begin{equation}\label{eq:beamweight}
 |D^q(\beta u_*)|\le C_qk^qW,\quad
 |D^q(L+k^2)(\beta u_*)|\le C_{q,M}k^{-M}W,\quad
 |\dd W|/W\le Ck^2 .
\end{equation}
Derivatives of the fixed collar cutoff can depend on \(w_{\rm c}\);
the positive gap absorbs them when \(k\) is sufficiently large.

\subsection{A real sample with a nonvanishing first jet}

Choose a net in \(\{t\ge s_1\}\) of mesh at most \(k^{-10}\).
There are at most \(Ck^{40}\) centers. At each take the four
packets of Proposition~\ref{prop:packets}, and denote the collection
by \(\mathcal W_k\). Use independent standard real Gaussian
coefficients \(g_v^{(1)},g_v^{(2)}\) and set
\begin{equation}\label{eq:random}
 u=\beta u_*+\sum_{v\in\mathcal W_k}
       (g_v^{(1)}\operatorname{Re}v+
                         g_v^{(2)}\operatorname{Im}v).
\end{equation}
Let \(\mathcal B_k\) be the event that all coefficients have
absolute value at most \(k\). The elementary Gaussian tail bound gives
\(\Prob(\mathcal B_k^c)\le Ck^{40}e^{-k^2/2}\to0\).

Given finite \(J\ge2\) and \(M>0\), choose the packet residual
orders through \(J\) to exceed \(M+42\). On \(\mathcal B_k\),
their number, \eqref{eq:packetder} and \eqref{eq:beamweight} give
\begin{equation}\label{eq:samplebounds}
 |D^q u|\le C_qk^{2q+41}W,\qquad
 |D^q\mathfrak r|\le C_qk^{-M}W\quad(0\le q\le J),\qquad
 \mathfrak r=\Div_\mu(A\,\dd u)+k^2\rho u.
\end{equation}
All orders are chosen before \(k\).
Every packet vanishes for
\(t<s_1-Ck^{-1/3-100\delta}\), so in that region \(u=u_*\)
and \(\mathfrak r=0\).

\begin{lemma}\label{lem:jet}
For every fixed choice of packet orders,
\[
 \Prob\left(\mathcal B_k\cap
   \left\{\inf_B\frac{|u|+|\dd u|_{g_0}}{W}\ge k^{-500}\right\}\right)
 \longrightarrow1 .
\]
\end{lemma}

\begin{proof}
Where \(W/e^{kh}>k^{50}\), the cutoff is one for large \(k\),
by \eqref{eq:gluegaps}. The cosine and angular derivative of
the uncut beam give
\(|u_*|+|\dd u_*|_{g_0}\ge ce^{kb_N}\).
On \(\mathcal B_k\) the packet first jets are at most
\(Ck^{43}e^{kh}\), so the beam dominates there and gives a
lower bound \(cW\). This region includes \(t<s_1\).

In the remaining region there is a nearby center. The real
jet spanning from Proposition~\ref{prop:packets}, converted
to ordinary derivatives, gives a covariance matrix for the
five-dimensional value--gradient vector whose least eigenvalue
has square root at least
\[
 ck^{-1}e^{kh}\ge ck^{-51}W\ge k^{-52}W.
\]
One packet family suffices; others add nonnegative covariance.
The deterministic beam shifts only the mean.
Hence the density of the normalized vector
\(\mathcal J=W^{-1}(u,\partial_1u,\ldots,\partial_4u)\)
is at most \(Ck^{260}\). For fixed \(C_0,D>0\),
\begin{equation}\label{eq:smallball}
 \Prob(|\mathcal J(x)|\le C_0k^{-D})\le Ck^{260-5D}.
\end{equation}
This is an unconditional Gaussian estimate.

On \(\mathcal B_k\), two derivatives in
\eqref{eq:samplebounds} and \(|\dd W|/W\le Ck^2\) give
\(|D\mathcal J|\le Ck^{45}\le k^{46}\).
In a finite atlas take grids of mesh \(k^{-D-48}\), with
at most \(Ck^{4D+192}\) points. Apply \eqref{eq:smallball}
at the Gaussian-region grid points and the beam bound elsewhere.
The union bound is \(Ck^{452-D}\). With \(D=500\) it tends
to zero, while the interpolation error is \(O(k^{-D-2})\).
Choose \(C_0\) to cover the fixed coordinate norm comparisons.
This yields the stated global first-jet bound, after also
subtracting \(\Prob(\mathcal B_k^c)\).
The use of \(\mathcal B_k\) is confined to derivative control
and interpolation, not to the Gaussian density calculation.
\end{proof}

\subsection{Many strict signs at the local wavelength}

\begin{lemma}\label{lem:signs}
For fixed packet orders there are \(p_*,c_*>0\) such that,
for all sufficiently large allowed \(k\), \(u\) has a sign
certificate of size at least \(c_*k^{1+\gamma}\) with
probability at least \(p_*\). Potential disks can be fixed
deterministically before sampling.
\end{lemma}

\begin{proof}
In \(V\), use radii \(R_x=(kQ_x)^{-1}\).
Since \(|D Q|\le Ck^{4\delta}\), intersecting balls have
comparable \(Q\)'s and radii for large \(k\).
At a fixed \(k\) these radii have a positive minimum.
Choose a finite maximal disjoint family
\(\{B(x,R_x):x\in\mathcal X_k\}\).
Fixed enlargements cover \(V\), and \(Q\) is comparable
to \(Q_x\) in each enlargement. Therefore
\begin{equation}\label{eq:packing}
 \int_V kQ\,\dd x\le C\sum_{x\in\mathcal X_k}
             kQ_xR_x^4=C\sum_xR_x^3,\qquad
 \sum_xR_x^3\ge ck^{1+\gamma}.
\end{equation}
The latter inequality uses the integral in \eqref{eq:envelope}.

The beam vanishes on radius-two enlargements of these balls.
On one define
\[
 S_x^2=\sum_v|v(x)|^2,\qquad
 U_x(Y)=u(x+R_xY)/S_x,\quad |Y|\le2.
\]
Lemma~\ref{lem:scaled} gives all fixed-order squared derivative
bounds after this normalization, and the non-tail value fraction
is \(1-o(1)\). The value \(U_x(0)\) is a standard real Gaussian.

Let \(z\) be the coefficient vector with pair
\((\operatorname{Re}v(x),\operatorname{Im}v(x))\) for each packet,
and \(T\) the matrix of their four real \(Y\)-derivatives at zero.
Then \(|z|=S_x\), and the gradient covariance conditional on
the value is
\begin{equation}\label{eq:condcov}
 S_x^{-2}T\left(I-\frac{zz^{\mathsf t}}{S_x^2}\right)T^{\mathsf t}.
\end{equation}
The formula follows by an orthogonal decomposition of the
standard Gaussian coefficient vector along \(z\).
For each non-tail, the unit coefficient direction
\[
 c_v=(-\operatorname{Im}v(x),\operatorname{Re}v(x))/|v(x)|
\]
is perpendicular to \(z\); these directions are mutually
orthonormal, and
\[
 T c_v=|v(x)|\operatorname{Im}\frac{\dd_Yv}{v}(0).
\]
Equations~\eqref{eq:imaginary} and \eqref{eq:squares} give a
fixed positive trace lower bound in \eqref{eq:condcov}.
There is a deterministic unit direction \(e_x\) such that
\(Z_x=\partial_{e_x}U_x(0)\), conditional on \(U_x(0)\),
has variance between fixed positive constants. Unconditional
variance and the regression coefficient are bounded above
by the derivative square sums.

Those same estimates through order five, integrated on
\(B(0,2)\subset\R^4\), give
\begin{equation}\label{eq:c2moment}
 \E\|U_x\|_{C^2(\overline{B(0,1)})}^2
 \le C\,\E\|U_x\|_{H^5(B(0,2))}^2\le C.
\end{equation}
The Sobolev inequality here is the \(k=5,m=2,p=2,n=4\)
case of \cite[Theorem~3.49, p.~76]{Hunter}.
In particular, it controls the Euclidean gradient and
Hessian operator norms.

Take \(a'=1/4\), \(0<c'<a'/4\). For each sufficiently large
fixed \(H\), the conditional Gaussian density gives
\[
 \Prob(1\le Z_x\le2,\ |U_x(0)|<c'/H)\ge c_1/H.
\]
Its conditional mean is bounded on the indicated value interval,
and its conditional variance stays above and below fixed
positive numbers. Markov's inequality applied to
\eqref{eq:c2moment} costs at most \(C/H^2\) for the
failure of \(\|U_x\|_{C^2}\le H\).
Fix \(H\) large enough so the intersection has probability
\(\pi_*=c_1/(2H)>0\).

Taylor's formula at \(\pm a'e_x/H\) now gives opposite
values of absolute size at least \(a'/(2H)\).
The gradient bound keeps those signs on the parallel transverse
disks of radius \(\sigma_*=a'/(4H^2)\).
Their joining cylinder lies inside the unit ball.
The directions, disks and radii are deterministic; only success
depends on the Gaussian sample.

Let \(F_k\) be the fraction of successful balls weighted by
\(R_x^3\). Then \(0\le F_k\le1\), \(\E F_k\ge\pi_*\).
Writing \(q=\Prob(F_k\ge\pi_*/2)\) gives
\(\pi_*\le(\pi_*/2)(1-q)+q\), hence
\[
 \Prob(F_k\ge\pi_*/2)\ge\pi_* /(2-\pi_*)>0.
\]
No independence between balls is needed.
Each disk has three-area proportional to
\(\sigma_*^3R_x^3\). Summing and using \eqref{eq:packing}
proves the lemma.
\end{proof}

After all finite orders are fixed, Lemma~\ref{lem:signs}
has one positive probability \(p_*\), while failure in
Lemma~\ref{lem:jet} tends to zero.
For every sufficiently large allowed \(k\), probability
subtraction therefore gives a real sample satisfying
\eqref{eq:samplebounds}, the global first-jet lower bound,
and a certificate of size \(c_*k^{1+\gamma}\).

\subsection{Canceling the residual exactly}

\begin{lemma}\label{lem:repair}
Suppose the first-jet bound in Lemma~\ref{lem:jet} holds.
For an integer \(R\ge0\), assume the derivative bounds
\eqref{eq:samplebounds} for \(u\) through \(R+2\) and for
\(\mathfrak r\) through \(R+1\). There are smooth real corrections
to the pair, zero on every open set where \(\mathfrak r=0\), such that
\[
 \Div_\mu((A+\delta A)\dd u)+k^2(\rho+\delta\rho)u=0,\qquad
 \|(\delta A,\delta\rho)\|_{C^R}
       \le C_Rk^{-M+2129+1088R}.
\]
Thus \(M\ge2131+1088R\) gives corrections tending to zero in \(C^R\).
\end{lemma}

\begin{proof}
Put \(I=g_0^{-1}\), and define
\begin{equation}\label{eq:repair}
 \mathcal D=u^2+I(\dd u,\dd u),\quad
 F=-\mathfrak r/\mathcal D,\quad
 \delta A=FuI,\quad
 \delta\rho=k^{-2}\bigl(Fu-\Div_\mu(FI\dd u)\bigr).
\end{equation}
The first-jet bound implies
\(\mathcal D\ge \tfrac12 k^{-1000}W^2>0\).
Consequently all four expressions are globally smooth. The product rule
gives
\[
 \Div_\mu(FuI\dd u)
   +\bigl(Fu-\Div_\mu(FI\dd u)\bigr)u
     =F\bigl(u^2+I(\dd u,\dd u)\bigr)=-\mathfrak r.
\]
This is the required exact identity.

For \(0\le q\le R+1\), the bounds on \(u\) give
\[
 |D^q\mathcal D|\le C_q k^{2q+86}W^2.
\]
A derivative of order \(q\ge1\) of \(\mathcal D^{-1}\) is a sum
of terms
\(\mathcal D^{-m-1}\prod_{\nu=1}^m D^{q_\nu}\mathcal D\),
where \(q_\nu\ge1\), \(\sum q_\nu=q\), and \(m\le q\).
The lower bound for \(\mathcal D\) therefore yields, also at order zero,
\[
 |D^q\mathcal D^{-1}|
    \le C_q k^{1000+1088q}W^{-2},\qquad
 |D^qF|\le C_q k^{-M+1000+1088q}W^{-1}.
\]
These are estimates on already differentiated factors; they do not
differentiate \(W^{-1}\).
Multiplication by derivatives of \(u\) cancels the weights and gives
\[
 |D^q(Fu)|\le C_qk^{-M+1041+1088q},\qquad
 |D^q(FI\dd u)|\le C_qk^{-M+1043+1088q}.
\]
The divergence in \eqref{eq:repair} costs one derivative, followed
by \(k^{-2}\). Its exponent is
\(1043+1088(R+1)-2=2129+1088R\), as stated.
If \(\mathfrak r\) vanishes on an open set, then \(F\) and its
derivatives vanish there, so both corrections vanish there.
For sufficiently large \(k\) they preserve positivity, since the
input pair is positive on a compact manifold.
\end{proof}

\subsection{Making the eigenvalue simple while retaining its eigenfunction}

\begin{lemma}\label{lem:simple}
Suppose \(u\) is an exact real eigenfunction at \(k^2=N(N+3)\),
with \(N>1\). Suppose the pair is round and \(u=u_*\) on an open
region containing a point where \(0<t<1\) and \(\dd u,\dd t\)
are independent. A smooth tensor perturbation supported there can
make \(k^2\) simple while retaining \(u\) and \(\rho\).
Its size can be arbitrarily small in any prescribed finite \(C^R\) norm.
\end{lemma}

\begin{proof}
Take a small product coordinate patch \(O\) with coordinates
\((u,t,z_3,z_4)\). Choose a nonnegative bump \(\chi\) supported
inside \(O\), positive on a smaller product patch \(O_0\). On \(O\) put
\[
 E=\chi\left(I-
   \frac{(I\dd u)\otimes(I\dd u)}{I(\dd u,\dd u)}\right),
 \qquad I=g_0^{-1},
\]
and extend it by zero. The tensor is smooth, nonnegative and
\(E\dd u=0\). Let \(\mathcal E\) be the old real eigenspace at
\(k^2\). We first show that the kernel on \(\mathcal E\) of
\(\int_B E(\dd f,\dd f)\dd\mu\) is exactly \(\R u\).

If the integral is zero, positivity on \(O_0\) gives
\(\dd f\in\R\,\dd u\) there. On its product coordinates this
means \(f=b(u)\), with \(b\) smooth. The round equations imply
\begin{equation}\label{eq:composition}
 b''(u)|\dd u|_{g_0}^2=k^2\bigl(ub'(u)-b(u)\bigr).
\end{equation}
The ambient gradient of \(\operatorname{Re}(x_1+i x_2)^N\)
has squared norm \(N^2(1-t)^{N-1}\) and radial derivative \(Nu\).
Projection to the tangent space of the sphere gives
\[
 |\dd u|_{g_0}^2=N^2\bigl((1-t)^{N-1}-u^2\bigr).
\]
Differentiate \eqref{eq:composition} in the coordinate \(t\),
holding \(u\) fixed. Since \(N>1\) and \(t<1\), this gives
\(b''=0\); substitution eliminates the constant term of \(b\).
Thus \(f=cu\) on \(O_0\).
The difference satisfies a smooth uniformly elliptic second-order
equation. Aronszajn's theorem \cite[pp.~235--236]{Aronszajn}
applies: the smooth principal coefficients satisfy its
\(C^2\) and Lipschitz-second-derivative requirements, and the
lower-order coefficients are bounded locally. Open-set vanishing
gives infinite-order vanishing at an interior point. Unique
continuation, applied in overlapping charts on connected \(B\),
therefore gives \(f=cu\) globally.

Now replace \(A\) by \(A+\tau E\), with \(\tau>0\).
The same \(u\) still solves the exact equation at \(k^2\).
If that eigenvalue were not simple for arbitrarily small positive
\(\tau\), there would be \(\tau_m\downarrow0\) and corresponding
real eigenfunctions \(f_m\), of unit weighted \(L^2\) norm and
orthogonal to \(u\). Their energy identity
\[
 \int_B(A+\tau_mE)(\dd f_m,\dd f_m)\dd\mu=k^2
\]
bounds them in \(H^1\). Weak \(H^1\) compactness and compact
\(H^1\hookrightarrow L^2\) give a subsequence converging to
a norm-one \(f_0\) orthogonal to \(u\).
The weak equation passes to the limit; local elliptic regularity
makes \(f_0\) smooth, so \(f_0\in\mathcal E\).
For any fixed \(\phi\in\mathcal E\), testing the old equation
against \(f_m\) and the perturbed one against \(\phi\), then
subtracting and dividing by \(\tau_m\), gives
\[
 \int_B E(\dd f_m,\dd\phi)\dd\mu=0.
\]
Weak convergence preserves this identity. Taking \(\phi=f_0\)
places \(f_0\) in the kernel already shown to be \(\R u\),
which contradicts its norm and orthogonality.
Thus every sufficiently small positive \(\tau\) makes \(k^2\)
simple. The tensor \(E\) is fixed before this scalar is chosen,
so \(\tau E\) can meet any finite derivative tolerance.
The compactness and regularity used here are justified explicitly
in Lemma~\ref{lem:persistence} below.
\end{proof}

\subsection{Completion of the quantitative step}

\begin{proof}[Proof of Proposition~\ref{prop:step}]
Fix the pair, collar, \(r,\eta\). Choose \(R\ge\max(r,1)\) and
a smaller positive tolerance \(\eta'\) such that movement
less than \(\eta'\) in \(C^R\) implies the desired \(C^r\)
bound and keeps the pair positive and inside \(\mathcal U\).
This is possible because the fixed input is in the open neighborhood.
Set \(J=R+2\), take \(M\ge2131+1088R\), and then choose the
finite packet orders for \eqref{eq:samplebounds}.
These choices precede the frequency.

For every sufficiently large allowed \(k\), the simultaneous
sample above has a certificate of size \(c_*k^{1+\gamma}\)
and the hypotheses of Lemma~\ref{lem:repair}.
Choose \(k\) sufficiently large that its correction is smaller
than \(\eta'/2\) in \(C^R\).
The protected round region remains nonempty for every fixed
collar width: it suffices to require
\begin{equation}\label{eq:collarthreshold}
 Ck^{-1/3-100\delta}<w_{\rm c}/100,\qquad
 c w_{\rm c}k>(M'+P)\log k+\log C_q(w_{\rm c}).
\end{equation}
The second inequality makes a cutoff error
\(C_q(w_{\rm c})k^P e^{-c w_{\rm c}k}\) smaller than
\(k^{-M'}\). There are only finitely many derivative orders
and accuracy targets at this stage, so all these inequalities
hold for every sufficiently large \(k\).
Below \(s_1-w_{\rm c}/100\), \(u=u_*\) and the residual and
corrections vanish.

In \(l_-+w_{\rm c}/20<t<l_-+w_{\rm c}/10\), choose a point
where \(\sin(N\vartheta)\ne0\). Here \(\dd t\ne0\) and
the angular derivative of \(u_*\) is nonzero, so \(\dd u,\dd t\)
are independent. Apply Lemma~\ref{lem:simple}. After \(k\)
and the sample have been fixed, choose its scalar \(\tau\)
below the simplicity threshold with
\(\|\tau E\|_{C^R}<\eta'/2\).
All derivatives of \(E\) are finite at this stage; their possible
growth with \(k\) causes no obstruction to this later choice.
The perturbation is supported in \(t>l_-\).
It preserves the function, eigenvalue and signs exactly.

The final pair is round for \(t<l_-\), meets the norm tolerance,
and has a simple eigenvalue and the required certificate.
The number \(c_*>0\) is fixed before the universal choice of
large frequency. The exponent \(\gamma\), chart \(\Omega\)
and neighborhood \(\mathcal U\) were fixed earlier, independent
of all the finite orders, tolerances and collar widths.
\end{proof}

\section{Persistence and one fixed smooth metric}\label{sec:limit}

\subsection{A finite norm suffices to preserve earlier signs}

\begin{lemma}\label{lem:persistence}
Let a smooth positive pair \(\mathcal P\) have a simple positive
eigenvalue \(\lambda\) and a real eigenfunction with a sign
certificate \(\mathcal C\). For any open interval
\(I\subset(0,\infty)\) containing \(\lambda\), there are
an integer \(s\ge1\) and \(\varepsilon>0\) such that every
smooth positive pair \(\mathcal P'\) with
\(\|\mathcal P'-\mathcal P\|_{C^s}\le\varepsilon\)
has a simple eigenvalue in \(I\) and a real eigenfunction
with the same certificate disks.
\end{lemma}

\begin{proof}
We specify the compactness and regularity facts used both here
and in Lemma~\ref{lem:simple}.
On \(H^1(B)\), the shifted form
\[
 a(f,v)=\int_B A(\dd f,\dd v)\dd\mu+
                     \int_B\rho fv\dd\mu
\]
is continuous, symmetric and coercive.
The inclusion \(H^1(B)\hookrightarrow L^2(\rho\dd\mu)\)
is continuous, compact and dense.
These properties follow in a finite atlas from the local Sobolev
compactness theorem \cite[Theorem~3.49, p.~76]{Hunter};
positive smooth densities give equivalent \(L^2\) norms.
Laugesen's discrete spectral theorem
\cite[Theorem~4.1, pp.~28--32]{Laugesen}, with form space
\(H^1\) and Hilbert space \(L^2(\rho\dd\mu)\), gives an orthonormal
eigenbasis for the shifted problem with eigenvalues tending to
infinity. Subtracting one gives the nonnegative discrete spectrum
of \(-L\), with finite multiplicities.
All eigenfunctions can be chosen real.

For completeness, weak eigenfunctions are smooth.
In a coordinate chart write \(\dd\mu=m(x)\dd x\).
The divergence coefficients are \(mA^{ab}\), and the right-hand
side is a smooth density times the eigenfunction.
Local uniform ellipticity and
\cite[Theorems~4.27--4.28, pp.~112--114]{Hunter}
first give \(H^2_{\rm loc}\) and then successively every finite
Sobolev order. Their hypotheses on coefficient derivatives and
the Sobolev regularity of the right-hand side are met at each
step. Sobolev embedding and a finite cover give global smoothness.
The same argument applies to a weak limit solving a fixed smooth
eigenfunction equation.

The eigenbasis also gives the min--max characterization
\begin{equation}\label{eq:minmax}
 \lambda_n=\inf_{\dim F=n+1}\ \sup_{0\ne f\in F}
 \frac{\int_B A(\dd f,\dd f)\dd\mu}
      {\int_B\rho f^2\dd\mu},\qquad n=0,1,\ldots .
\end{equation}
Indeed, the first \(n+1\) eigenvectors give the upper bound.
Every \((n+1)\)-dimensional \(F\) contains a nonzero vector
orthogonal to the first \(n\), whose eigenbasis expansion gives
the lower bound. The characterization is on the common form
domain \(H^1\), including when the density varies.
Uniform multiplicative comparisons of nearby positive tensors
and densities compare all the quotients in \eqref{eq:minmax}.
Thus \(C^0\) convergence of pairs implies convergence of every
fixed ordered eigenvalue. A smaller interval with closure in \(I\)
can separate the simple \(\lambda\) from its neighboring
ordered eigenvalues. For all sufficiently close pairs it contains
exactly the corresponding eigenvalue, which is still simple.

We need convergence of its normalized eigenfunction in \(C^0\).
The same interior estimates, localization by a finite smooth
partition of unity and absorption of lower-order terms give,
for the fixed smooth operator \(L\),
\begin{equation}\label{eq:elliptic}
 \|f\|_{H^{a+2}}\le C_a
       \bigl(\|Lf\|_{H^a}+\|f\|_{L^2}\bigr),\qquad a\ge0.
\end{equation}
One way to see the absorption is to apply the interior estimates
on a slightly larger chart to cutoff multiples of \(f\).
The commutator with the cutoff is first order.
The inequalities
\(\|f\|_{H^{a+1}}\le
 \epsilon\|f\|_{H^{a+2}}+C_\epsilon\|f\|_{L^2}\)
follow from the elementary Fourier weight inequality in each
chart. A finite sum and sufficiently small \(\epsilon\) absorb
those commutators and give \eqref{eq:elliptic}.

A sufficiently small \(C^4\) change of the pair makes
\(L'-L\) small as a map \(H^{a+2}\to H^a\) for \(a=0,2\).
In coordinates its principal coefficients have no coefficient
derivatives, and its first-order coefficients have at most one;
the stated finite norm bounds every product-rule coefficient
needed at these two orders. Absorb
\(\|(L'-L)f\|_{H^a}\) into \eqref{eq:elliptic}.
For eigenvalues in the chosen bounded interval, the estimates
at \(a=0\) and then \(a=2\) give uniform \(H^4\) bounds for
weighted-\(L^2\)-normalized eigenfunctions.

If \(\mathcal P_m\to\mathcal P\) in this finite norm, choose
real eigenfunctions \(f_m\) with
\(\int_B\rho_m f_m^2\dd\mu=1\).
A subsequence converges strongly in \(H^3\), and hence in \(C^0\);
the dimension-four embedding \(H^3\hookrightarrow C^0\)
again follows from Hunter's Theorem~3.49 in a finite atlas.
The weak equation passes to the limit, whose norm is one
because \(\rho_m\to\rho\) uniformly. Simplicity identifies the
limit with the given normalized eigenfunction, up to sign.
On the finite collection of compact disks its strict signs
have a positive absolute margin. Uniform convergence retains
all those signs, possibly after a global sign change.
If no sufficiently small finite-norm neighborhood retained
them, a sequence contradicting this conclusion could be
chosen and the same compactness argument would apply.
Finally shrink the radius to put a closed finite-norm ball
inside that neighborhood.
\end{proof}

\subsection{Summable changes and separated exact eigenvalues}

\begin{proof}[Proof of Theorem~\ref{thm:main}]
Take monotone coordinate norms as in Section~\ref{sec:base}.
Choose \(b>0\) such that the closed \(C^1\) ball of radius \(b\)
about \(\mathcal P_0=(g_0^{-1},1)\) consists of positive pairs
inside \(\mathcal U\). Set
\[
 l^{(j)}=\frac14+2^{-j-2}\quad(j\ge0),\qquad
 \epsilon_0=\gamma/4>0 .
\]
We construct \(\mathcal P_j=(A_j,\rho_j)\), round for
\(t<l^{(j)}\), with a simple eigenvalue \(k_j^2\) and
certificate \(\mathcal C_j\).

After stage \(i\), apply Lemma~\ref{lem:persistence} with interval
\((k_i^2/2,2k_i^2)\), and call its norm order and radius
\(s_i,\varepsilon_i\). At stage \(j\), only earlier such
numbers enter the requirements:
\begin{equation}\label{eq:budgets}
 \begin{split}
 r_j&=\max\{j+1,1,s_1,\ldots,s_{j-1}\},\\
 \eta_j&=\min\left\{2^{-j},\,b2^{-j-1},\,
              \min_{1\le i<j}
                    2^{-(j-i+1)}\varepsilon_i\right\}.
 \end{split}
\end{equation}
For \(j=1\) omit the inner minimum. Apply
Proposition~\ref{prop:step} to the fixed input
\(\mathcal P_{j-1}\), collar
\((l^{(j)},l^{(j-1)})\), order \(r_j\) and tolerance \(\eta_j\).
Its positive constant \(c_j\) is known before choosing \(k_j\).
Choose an allowed frequency beyond its threshold satisfying
\[
 k_j>2k_{j-1}\quad(j>1),\qquad c_jk_j^{\gamma/2}>j.
\]
The resulting certificate then obeys
\begin{equation}\label{eq:stagesize}
 |\mathcal C_j|>j k_j^{1+\gamma/2}.
\end{equation}
Use the disks constructed in Lemma~\ref{lem:signs}; their centers
are in \(V\Subset\Omega\), and their radii are at most \(k_j^{-1}\).
Choose also the first frequency large enough to keep all cylinders
in one fixed compact subregion of \(\Omega\).
Only now select \(s_j,\varepsilon_j\) for this output.
They restrict later stages and do not constrain the frequency
that produced them.

For every fixed \(r\), eventually \(r_j\ge r\), and
\(\|\mathcal P_j-\mathcal P_{j-1}\|_{C^r}<2^{-j}\).
The sequence therefore converges in every \(C^r\) to a single
smooth pair \(\mathcal P_\infty=(A_\infty,\rho_\infty)\).
Its total \(C^1\) movement from the round pair is at most
\(b\sum_{j\ge1}2^{-j-1}=b/2\), so it is positive and lies
inside \(\mathcal U\).
For each \(i\), \eqref{eq:budgets} gives
\[
 \|\mathcal P_\infty-\mathcal P_i\|_{C^{s_i}}
 \le\sum_{j>i}2^{-(j-i+1)}\varepsilon_i
       =\varepsilon_i/2 .
\]
Thus Lemma~\ref{lem:persistence} applies simultaneously to every
earlier certificate. There are nonzero real smooth \(v_i\) for
the same limiting operator satisfying
\begin{equation}\label{eq:limitlambda}
 -L_\infty v_i=\lambda_i v_i,\qquad
 k_i^2/2<\lambda_i<2k_i^2,
\end{equation}
with certificate disks \(\mathcal C_i\).
The eigenvalue intervals are disjoint and tend to infinity
because \(k_{i+1}>2k_i\).

\subsection*{An ordinary metric with the prescribed base operator}

Put
\[
 G_\infty^{-1}=A_\infty/\rho_\infty,\qquad
 q=\frac{\rho_\infty\,\dd\mu}{\dd\operatorname{vol}_{G_\infty}},
 \qquad
 g=G_\infty+q^2\dd\varphi^2\quad\hbox{on }B\times S^1 .
\]
The ratio \(q\) of smooth positive densities is smooth and positive.
Hence \(g\) is a smooth positive metric on a closed connected
five-dimensional manifold, fixed independently of \(i\).
Its volume density is
\[
 \dd\operatorname{vol}_g
      =q\,\dd\operatorname{vol}_{G_\infty}\dd\varphi
      =\rho_\infty\,\dd\mu\dd\varphi .
\]
For any function \(f\) independent of \(\varphi\), the coordinate
divergence formula, with \(\mu=m(x)\dd x\), gives
\begin{equation}\label{eq:circle}
 \Delta_g f=\frac{1}{\rho_\infty m}
  \partial_a(\rho_\infty mG_\infty^{ab}\partial_b f)
     =\rho_\infty^{-1}\Div_\mu(A_\infty\dd f)=L_\infty f .
\end{equation}
The pullbacks \(u_i\) of \(v_i\) are therefore exact real
Laplace eigenfunctions on this single metric, at \(\lambda_i\).

\subsection*{An intrinsic Hausdorff measure lower bound}

Fix a closed interval \(I_0\) of positive length in a circle
coordinate chart. For one certificate cylinder \(C_i^{(a)}\),
let \(D_i^{(a)}\) denote its transverse disk.
Continuity and opposite endpoint signs force a zero on each
joining segment. Consequently
\[
 \Z(u_i)\cap(C_i^{(a)}\times I_0)
\]
projects onto \(D_i^{(a)}\times I_0\) by the orthogonal
base-coordinate projection and the circle coordinate.

All these cylinders lie in a single compact coordinate region.
Extend its coordinate functions, and the circle coordinate on
\(I_0\), smoothly with cutoffs to the compact manifold.
Their gradients have fixed bounds in \(g\). Integrating along
curves bounds their differences by a fixed constant times
the intrinsic distance. Composing with any orthogonal projection
gives one Lipschitz constant \(K\), uniform in every disk's
radius and orientation.
The Hausdorff covering inequality therefore implies
\begin{equation}\label{eq:project}
 \HH_g^4\bigl(\Z(u_i)\cap(C_i^{(a)}\times I_0)\bigr)
 \ge K^{-4}\HH_{\mathrm{Euc}}^4(D_i^{(a)}\times I_0)
 \ge c\,\HH_{\mathrm{Euc}}^3(D_i^{(a)})|I_0|.
\end{equation}
The last constant is purely dimensional: the product disk has
four-dimensional Lebesgue volume equal to the disk's
three-volume times \(|I_0|\), while a set of diameter \(d\)
in its four-dimensional affine span has volume at most \(Cd^4\).
Apply this estimate to every countable cover and take the
Hausdorff infimum. The identical three-dimensional covering
normalization identifies the disk's three-volume with a fixed
multiple of \(\HH^3_{\mathrm{Euc}}\).
Thus no convention on the normalization of Hausdorff measure
affects the bound.

The finitely many closed cylinders are in disjoint open balls,
so their compact nodal subsets are pairwise positively separated.
Hausdorff measure is additive on these separated sets directly
from covers of sufficiently small diameter. Summing
\eqref{eq:project} and using \eqref{eq:stagesize} yields
\begin{equation}\label{eq:finalmeasure}
 \HH_g^4(\Z(u_i))\ge c|\mathcal C_i|
                   >c\,i k_i^{1+\gamma/2}.
\end{equation}
Only continuity and strict signs were used; no regularity of
the nodal set is needed.
Finally, \eqref{eq:limitlambda} and \(\epsilon_0=\gamma/4\) give
\[
 \lambda_i^{1/2+\epsilon_0}
 \le 2^{1/2+\epsilon_0}k_i^{1+\gamma/2},\qquad
 \frac{\HH_g^4(\Z(u_i))}{\lambda_i^{1/2+\epsilon_0}}
       \ge c\,2^{-1/2-\epsilon_0}i\longrightarrow\infty .
\]
The metric and positive exponent are fixed for the full sequence.
This proves Theorem~\ref{thm:main}.
\end{proof}


\paragraph{Attribution and AI use.}
The counterexample, envelope construction, local waves, coefficient
correction and iteration are due to OpenAI's preprint \cite{OpenAI2026}.
This note adapts that publicly released construction into an exposition
of its implication for AIM Problem~011. Codex assisted with primary-source
retrieval, proof reconstruction, mathematical checking, drafting and
typesetting. These AI checks do not constitute independent human
certification. No formal proof build is claimed, and no numerical
experiment is used as evidence for the theorem.

\begin{thebibliography}{9}
\bibitem{OpenAI2026}
OpenAI,
\emph{Power-law violations of Yau's nodal upper bound},
OpenAI Math Release preprint, September 23, 2026.
Source revision \texttt{fd4aeeb2ee4fc729c18d98444fed42fd0529eeeb}.
\href{https://github.com/openai/math/tree/fd4aeeb2ee4fc729c18d98444fed42fd0529eeeb/preprints/Power-law-violations-of-Yaus-nodal-upper-bound-September-23-2026}{Pinned paper and editable source}.

\bibitem{AIM011}
AIM problem repository,
\emph{011. Yau's sharp upper bound for nodal volume},
revision \texttt{c9929805cece044705abaaebe1cebddac1ba11a6}.
\href{https://github.com/MColbrook/AIM/blob/c9929805cece044705abaaebe1cebddac1ba11a6/problems/011-yau-nodal-upper-bound.md}{Immutable problem statement}.

\bibitem{Aronszajn}
N. Aronszajn,
\emph{A unique continuation theorem for solutions of elliptic partial
differential equations or inequalities of second order},
Journal de Math\'ematiques Pures et Appliqu\'ees (9) \textbf{36}
(1957), 235--249.
\href{https://sites.math.washington.edu/~blwilson/Nodal/Aronszajn.pdf}{Original scan}.

\bibitem{Hunter}
John K. Hunter,
\emph{Notes on Partial Differential Equations},
University of California at Davis, revised June 18, 2014.
\href{https://www.math.ucdavis.edu/~hunter/pdes/pde_notes.pdf}{Author-hosted notes}.

\bibitem{Laugesen}
Richard S. Laugesen,
\emph{Spectral Theory of Partial Differential Equations---Lecture Notes},
arXiv:1203.2344v1 (2012).
The consulted PDF's title page is dated November 27, 2024.
\href{https://arxiv.org/pdf/1203.2344}{Lecture notes}.
\end{thebibliography}
\end{document}
